% Generated draft: requires mathematical and source review.
\documentclass[11pt]{article}
\usepackage[margin=25mm]{geometry}
\usepackage{fontspec}
\setmainfont{Latin Modern Roman}
\usepackage{amsmath,amssymb,mathtools,unicode-math}
\setmathfont{Latin Modern Math}
\usepackage{enumitem,needspace,xurl,hyperref}
\setlength{\parindent}{0pt}
\setlength{\parskip}{5pt}
\setlength{\emergencystretch}{3em}
\newcommand{\N}{\mathbb{N}}
\newcommand{\Z}{\mathbb{Z}}
\newcommand{\Q}{\mathbb{Q}}
\newcommand{\R}{\mathbb{R}}
\newcommand{\C}{\mathbb{C}}
\newcommand{\F}{\mathbb{F}}
\newcommand{\A}{\mathbb{A}}
\newcommand{\PP}{\mathbb{P}}
\newcommand{\E}{\mathbb{E}}
\newcommand{\eps}{\varepsilon}
\newcommand{\dd}{\,\mathrm{d}}
\begin{document}
\section*{UnsolvedMath ID: 30000660}
\textbf{Source problem code:} catalogue entry
\subsection{Definitions and mathematical statement}
\textbf{Fiberwise-isomorphism questions.} Let \(p:S\to X\) and \(q:T\to X\) be morphisms of varieties over an algebraically closed field \(k\), and suppose every schematic fiber \(S_x\) is isomorphic to \(T_x\). Must there be a dominant \'etale morphism \(U\to X\) with\[S\times_XU\cong T\times_XU\quad\text{as }U\text{-schemes}?\] If not, can counterexamples be given over fields such as \(\F_p\) or \(\Q\)? Both questions are preserved, but the latter requires a separate field and fiber convention because those fields are not algebraically closed.\par For a scheme point \(x\in X\), the schematic fiber is \(S_x=S\times_X\operatorname{Spec}\kappa(x)\), and an isomorphism with \(T_x\) should be over \(\kappa(x)\). Retaining the scheme structure is stronger than comparing underlying point sets. The source does not say whether \(x\) ranges over closed \(k\)-points, every scheme point including a generic point, or geometric points after extending residue fields. These hypotheses are different. In particular, when \(X\) is integral, including the generic point with an actual \(\kappa(\eta)\)-isomorphism already supplies a generic isomorphism, which for these finite-type families spreads out over a dense open subset. That open inclusion is itself dominant and \'etale. Thus the all-scheme-points reading and the closed-fiber reading cannot be treated as the same nontrivial problem. A geometric generic isomorphism is also different from an isomorphism over the original function field.\par Dominant means dense image, not necessarily surjective; a dominant \'etale morphism may omit a closed subset. The conclusion as written does not assert an \'etale cover of every point of \(X\), nor require the morphism to be finite. Its family isomorphism must commute with the maps to \(U\). None of these conclusions is strengthened silently. The source supplies no flatness, properness, smoothness, or automorphism-group hypothesis on the families.\par For a variant over \(\F_p\) or \(\Q\), one must state whether the fiber condition is over rational points, all residue fields, or their algebraic closures, and whether a counterexample is required to remain one after extending the constant field. The source does not fix those choices. The supplied triage checked 2026-08-21 labels the extracted assertion open but gives no precise resolution or theorem comparison. Its dated status claim does not decide these formulation ambiguities and requires independent verification.
\subsection{Short English statement}
Ask whether fiberwise isomorphisms become one family isomorphism after a dominant \'etale base change, and ask for arithmetic counterexamples if not. Closed, generic, and geometric fibers impose different hypotheses. Dominance also differs from an \'etale cover of all the base.
\subsection{Sources}
\begin{itemize}\item Local UnsolvedMath registry, permanent ID 30000660, code catalogue entry; \emph{\'Etale-Local Equivalence of Fiberwise Isomorphic Families}, Oberwolfach Reports Open Problems.\item \emph{Affine Algebraic Geometry (2007)}, Oberwolfach Reports, DOI 10.4171/owr/2007/01, pp. 5--82. \url{https://doi.org/10.4171/owr/2007/01}. Proposed year 2007; no proposer named.\item The same full report is the sole supplied literature source; no precise family-isomorphism or counterexample theorem is supplied.\item Literature triage checked 2026-08-21, status \texttt{open}; record updated 2026-09-16. The \texttt{corrected\_verified} field records notation normalization and preservation of both questions. Verify dated claims independently.\end{itemize}
\end{document}
